\section{Related Work}

\subsection{Benchmark Lifecycle and Displacement Dynamics}

The lifecycle of ML benchmarks has received increasing attention as the community grapples with benchmark saturation and the challenge of meaningful progress measurement. Ott et al.\ \cite{ott2022benchmark} analyze 3,765 benchmarks, finding that benchmark \emph{breadth} (versatility across tasks) correlates with longevity at the population level. This population-level correlation motivates our investigation of whether individual-benchmark breadth predicts displacement timing --- a finer-grained survival analysis question that Ott et al.\ do not address.

Koch et al.\ \cite{koch2021reduced} establish a 87-task taxonomy for Papers With Code and measure \emph{concentration} dynamics in benchmark adoption. Their methodology provides the h-e2 panel that underlies our analysis. Where Koch et al.\ characterize concentration trends across tasks, we use the same panel to test individual predictors of displacement events.

Paullada et al.\ \cite{paullada2021data} provide a qualitative governance framework for dataset lifecycles, arguing that benchmark datasets are sociotechnical artifacts whose trajectories are shaped by community norms, institutional pressures, and competitive incentives. Our work operationalizes their qualitative observations into a quantitative survival analysis framework.

\subsection{Benchmark Overuse and Community Replacement}

The ICLR 2025 workshop on benchmarking \cite{iclr2025benchmarking} explicitly foregrounds benchmark overuse as a driver of community-initiated replacement. Under this framing (our H2 saturation mechanism), a benchmark that has been widely adopted signals consensus that the task is ``solved,'' incentivizing the community to propose harder successors. The GLUE $\to$ SuperGLUE transition is a canonical example: broad community adoption of GLUE \cite{wang2018glue} was followed by recognition of GLUE's ceiling, driving rapid adoption of SuperGLUE \cite{wang2019superglue}.

Our null result partially challenges this framing at the individual benchmark level: even if overuse drives replacement in high-profile cases, submission count breadth at introduction year does not predict displacement timing across the broader population of 87 tasks.

\subsection{Community Adoption and Lock-In Theory}

Rogers' Diffusion of Innovations \cite{rogers2003diffusion} provides the theoretical foundation for our H1 lock-in mechanism: early adopters become opinion leaders whose investment in a technology creates network effects that slow transitions to alternatives. Our null result suggests that this investment, at the aggregate level measured by submission count, does not translate into reduced displacement hazard.

\subsection{Survival Analysis in Bibliometrics}

We use \texttt{CoxPHFitter} from the lifelines library \cite{davidson2019lifelines} for all Cox regressions. The L2 penalizer (penalizer = 0.1) achieves clean convergence, validated across 14/14 unit tests on the h-e2 panel. The key methodological challenge --- identifying time-independent predictors that do not collapse into temporal proxies --- is directly addressed by our FAIL FAST protocol.
